\section{Spectral packet coordinates, the Connes--Consani Laplacian, and the Scaling Site}
\label{sec:spectral-packet-coordinates}

The packet divisor of Chapters~15--17 has two further published spectral
realizations.  Connes and Consani first factor the multiplicative Fourier
transform of a scale-invariant Riemann sum by
\(\zeta(1/2-\mathrm iz)\), then realize the values
\((\rho-1/2)^2-1/4\) as the spectrum of
\(\Delta=H(1+H)\), and finally identify a quotient sheaf on the Scaling Site
with a Schwartz quotient that detects the critical zeros and their possible
multiplicities
\cite[Eq.~(2.2), Prop.~4.2, and the theorem on pp.~96--98]{connesConsani2023Hochschild}.
This chapter constructs the exact maps from the packet coordinates already
proved here to those three realizations.  In particular, the affine map is
restricted to its correct open strip, the Laplacian's two-to-one fibre is
displayed, and an entire four-point packet is represented by a real polynomial
whose roots are a conjugate pair of Laplacian eigenvalues.

\subsection{The integral and complex split-zero models}

The local source for this programme uses \(G(\mathbb Z)\), whereas
Chapters~14--17 use \(G(\mathbb C)\).  These are related by a literal
embedding; they are not silently identified.

For a nonzero commutative ring \(R\), recall
\[
 G(R)=R\sqcup\{\tau\},
\]
with
\[
 r\oplus r'=r+r',\qquad r\oplus\tau=\tau\oplus r=r,
\]
and
\[
 r\odot r'=rr',\qquad r\odot\tau=\tau\odot r=\tau,
 \qquad \tau\odot\tau=\tau.
\]
The additive zero of \(G(R)\) is \(\tau\), while the supported element over
the ring zero is \(e_R:=0_R\).  Put
\(A_R=\{\tau,e_R\}\).

\begin{proposition}[Exact change of coefficient ring]
\label{prop:spc-integral-complex-embedding}
The map
\begin{equation}
\label{eq:spc-semiring-embedding}
 j:G(\mathbb Z)\longrightarrow G(\mathbb C),
 \qquad
 j(n)=n+0\mathrm i,\qquad j(\tau)=\tau,
\end{equation}
is an injective unital semiring homomorphism.  Its image is
\(\mathbb Z\sqcup\{\tau\}\), and every fibre over a point of its image is a
singleton.  It intertwines multiplication by \(-1\), and its restriction
\[
 j|_{A_{\mathbb Z}}:A_{\mathbb Z}\overset{\sim}{\longrightarrow}A_{\mathbb C}
\]
is the Boolean-semiring isomorphism \(\tau\mapsto\tau\),
\(e_{\mathbb Z}\mapsto e_{\mathbb C}\).  Equivalently,
\[
 j\bigl(\operatorname{Fix}_{G(\mathbb Z)}(-1)\bigr)
 =
 \operatorname{Fix}_{G(\mathbb C)}(-1)=A_{\mathbb C}.
\]
The map \(j\) is not surjective.  The inclusions
\(A_R\hookrightarrow G(R)\) are not unit-preserving, because the internal
unit of \(A_R\) is \(e_R\), whereas the ambient unit is \(1_R\); this does
not alter the fact that \(j\) itself and \(j|_{A_{\mathbb Z}}\) are unital in
their stated categories.
\end{proposition}

\begin{proof}
For supported integers \(m,n\), preservation of the two operations is the
ordinary inclusion calculation
\[
 j(m\oplus n)=j(m+n)=j(m)+j(n),
 \qquad
 j(m\odot n)=j(mn)=j(m)j(n).
\]
If either entry is \(\tau\), the defining tables on both sides give the same
answer: \(\tau\) is the additive identity and multiplicative absorber.  Also
\(j(\tau)=\tau\) and \(j(1_{\mathbb Z})=1_{\mathbb C}\), so the map preserves
zero and one in the ambient semirings.  Injectivity and the image formula
follow from the injectivity of \(\mathbb Z\hookrightarrow\mathbb C\) together
with the fact that the adjoined symbol \(\tau\) is disjoint from both supported
copies.

For every \(x\in G(\mathbb Z)\),
\(j((-1)\odot x)=(-1)\odot j(x)\).  The fixed-set computation of
\cref{prop:ds-fixed-pair} is valid over either torsion-free ring:
\(-z=z\) implies \(2z=0\), hence \(z=0\), while \(\tau\) is fixed because it
is absorbing.  Thus the fixed sets are precisely the two displayed Boolean
pairs, and the restriction of \(j\) is a bijection respecting their complete
operation tables.  A complex number outside \(\mathbb Z\) is not in the image,
so \(j\) is not surjective.  Finally, \(e_R\ne1_R\) for the nonzero rings in
question, which proves the assertion about the two inclusions.
\end{proof}

\subsection{The affine spectral strip and all four symmetries}

Let
\[
 \mathfrak S_s=\{s\in\mathbb C:0<\Re s<1\},
 \qquad
 \mathfrak S_z=\{z\in\mathbb C:|\Im z|<1/2\},
\]
and retain the unscaled affine coordinate
\begin{equation}
\label{eq:spc-affine-coordinate}
 \Phi(s)=\mathrm i\left(s-\frac12\right).
\end{equation}

\begin{proposition}[Biholomorphic strip coordinate and equivariance]
\label{prop:spc-affine-strip}
The map \(\Phi:\mathfrak S_s\to\mathfrak S_z\) is a biholomorphism with
\begin{equation}
\label{eq:spc-affine-inverse}
 \Phi^{-1}(z)=\frac12-\mathrm iz.
\end{equation}
Writing \(s=\beta+\mathrm i\gamma\), one has the coordinate formula
\begin{equation}
\label{eq:spc-affine-real-coordinates}
 \Phi(s)=-\gamma+\mathrm i\left(\beta-\frac12\right).
\end{equation}
For the functional involution \(w(s)=1-s\) and conjugation
\(c(s)=\overline s\), put
\[
 \sigma(z)=-z,
 \qquad
 \kappa(z)=-\overline z.
\]
Then
\begin{equation}
\label{eq:spc-affine-equivariance}
 \Phi\circ w=\sigma\circ\Phi,
 \qquad
 \Phi\circ c=\kappa\circ\Phi.
\end{equation}
Consequently the packet
\(\{s,1-s,\overline s,1-\overline s\}\) maps, with repetitions removed, to
\begin{equation}
\label{eq:spc-spectral-packet}
 \{z,-z,-\overline z,\overline z\},
 \qquad z=\Phi(s).
\end{equation}
Moreover, \(\Re s=1/2\) if and only if \(\Phi(s)\in\mathbb R\).
\end{proposition}

\begin{proof}
The affine map has nonzero derivative \(\mathrm i\), and direct substitution
gives \eqref{eq:spc-affine-inverse}.  Formula
\eqref{eq:spc-affine-real-coordinates} follows from
\[
 \mathrm i\left(\beta-\frac12+\mathrm i\gamma\right)
 =-\gamma+\mathrm i\left(\beta-\frac12\right).
\]
It also shows
\[
 0<\beta<1
 \quad\Longleftrightarrow\quad
 -\frac12<\Im\Phi(s)<\frac12,
\]
which proves the strip bijection.  The two equivariance identities are the
literal calculations
\[
 \Phi(1-s)=\mathrm i(1/2-s)=-\Phi(s)
\]
and
\[
 \Phi(\overline s)
 =\mathrm i(\overline s-1/2)
 =-\overline{\mathrm i(s-1/2)}.
\]
Applying these generators and their composite gives the four displayed
spectral points.  Finally,
\(\Im\Phi(s)=\beta-1/2\), proving the last equivalence.
\end{proof}

\subsection{The common-zero divisor, including multiplicities}

For an even Schwartz function \(f\), set
\begin{equation}
\label{eq:spc-E-and-psi}
 (\mathcal Ef)(u)=u^{1/2}\sum_{n\geq1}f(nu),
 \qquad
 \psi_f(z)=\int_{\mathbb R_+^*}f(u)u^{1/2-\mathrm iz}\,d^*u,
\end{equation}
and let
\[
 \mathcal S(\mathbb R)^{\mathrm{ev}}_1
 =
 \left\{f\in\mathcal S(\mathbb R)^{\mathrm{ev}}:
          \int_{\mathbb R}f(x)\,dx=0\right\}.
\]
Connes and Consani prove that, for every such \(f\), the multiplicative
Fourier transform
\[
 T_f(z):=\mathcal F(\mathcal Ef)(z)
\]
satisfies
\begin{equation}
\label{eq:spc-CC-factorization}
 T_f(z)=g(z)\psi_f(z),
 \qquad
 g(z):=\zeta\left(\frac12-\mathrm iz\right),
\end{equation}
throughout \(\Im z>-1/2\)
\cite[Eq.~(2.2) and the continuation immediately following it]{connesConsani2023Hochschild}.

For a family of holomorphic functions \(\mathcal T\) on a domain, define its
common-zero divisor by
\begin{equation}
\label{eq:spc-common-divisor-definition}
 D_{\mathrm{com}}(\mathcal T)
 =
 \sum_z\left(\min_{T\in\mathcal T}\operatorname{ord}_zT\right)[z],
\end{equation}
where a point with minimum zero contributes nothing.  The identically zero
member is assigned the convention \(\operatorname{ord}_z(0)=+\infty\), and
the minimum is finite because of the explicit nonzero member used in the
proof.

\begin{theorem}[Exact common divisor in the open spectral strip]
\label{thm:spc-common-zero-divisor}
Let \(D_\xi\) be the completed-zeta zero divisor, including multiplicities,
and let
\[
 \mathcal T_E=\{T_f:f\in\mathcal S(\mathbb R)^{\mathrm{ev}}_1\}.
\]
Then, as effective divisors on \(\mathfrak S_z\),
\begin{equation}
\label{eq:spc-three-divisors}
 \Phi_*D_\xi
 =
 \operatorname{div}_0(g)|_{\mathfrak S_z}
 =
 D_{\mathrm{com}}(\mathcal T_E)|_{\mathfrak S_z}.
\end{equation}
Thus both the support and the vanishing order at every point agree.  The
inverse from the common divisor to the completed divisor is the affine map
\(z\mapsto1/2-\mathrm iz\), and every point fibre of \(\Phi\) is a singleton.
\end{theorem}

\begin{proof}
All zeros of \(\xi\) lie in \(\mathfrak S_s\), and at such a zero \(\rho\)
the explicit factor
\[
 \frac12s(s-1)\pi^{-s/2}\Gamma(s/2)
\]
is holomorphic and nonzero.  Hence \(\xi\) and \(\zeta\) have the same
vanishing order at \(\rho\).  Since
\[
 g(z)=\zeta(\Phi^{-1}(z))
\]
and \((\Phi^{-1})'(z)=-\mathrm i\ne0\), composition with
\(\Phi^{-1}\) preserves that order.  This proves the first equality in
\eqref{eq:spc-three-divisors}, point by point and with multiplicity.

Equation \eqref{eq:spc-CC-factorization} implies
\[
 \operatorname{ord}_{z_0}T_f
 \geq
 \operatorname{ord}_{z_0}g
\]
for every \(f\) and every \(z_0\in\mathfrak S_z\).  It remains to prove that
the inequality is attained simultaneously at every point.  Use the explicit
even Schwartz function from Connes--Consani,
\begin{equation}
\label{eq:spc-explicit-witness}
 f_0(x)=e^{-\pi x^2}(2\pi x^2-1).
\end{equation}
The Gaussian identities
\[
 \int_{\mathbb R}e^{-\pi x^2}\,dx=1,
 \qquad
 \int_{\mathbb R}x^2e^{-\pi x^2}\,dx=\frac1{2\pi}
\]
give \(\int f_0=0\), so \(f_0\in\mathcal S(\mathbb R)^{\mathrm{ev}}_1\).
The cited computation is
\begin{equation}
\label{eq:spc-explicit-psi}
 \psi_{f_0}(z)
 =
 \frac14\pi^{-1/4+\mathrm iz/2}
 (-1-2\mathrm iz)
 \Gamma\left(\frac14-\frac{\mathrm iz}{2}\right).
\end{equation}
Each factor is finite and nonzero in \(\mathfrak S_z\).  The exponential
power of \(\pi\) never vanishes.  The linear factor vanishes only at
\(z=\mathrm i/2\), which is on the boundary.  If \(|\Im z|<1/2\), then
\[
 0<\Re\left(\frac14-\frac{\mathrm iz}{2}\right)<\frac12;
\]
hence its argument is not a pole of \(\Gamma\), and \(\Gamma\) has no zeros.
Therefore \(\psi_{f_0}\) is holomorphic and zero-free throughout the open
strip.  It follows that
\[
 \operatorname{ord}_{z_0}T_{f_0}
 =
 \operatorname{ord}_{z_0}g
\]
for every \(z_0\in\mathfrak S_z\), proving the second divisor equality.
The inverse and fibre assertions follow from
\eqref{eq:spc-affine-inverse}.
\end{proof}

\begin{remark}[Why the strip cannot be deleted]
\label{rem:spc-strip-essential}
The first equality in \eqref{eq:spc-three-divisors} is not a global equality
of zero sets on \(\mathbb C\).  For every \(n\geq1\), the trivial zero
\(s=-2n\) gives
\[
 z=\Phi(-2n)=-\mathrm i\left(2n+\frac12\right),
 \qquad
 g(z)=0,
\]
which lies below \(\Im z=-1/2\) and is not the image of a zero of \(\xi\).
The pole \(s=1\) of \(\zeta\) maps to the other boundary point
\(z=\mathrm i/2\).  The correct statement is therefore exactly the
open-strip divisor identity proved above.
\end{remark}

\subsection{The Laplacian map, its quotient, and its exceptional point}

Connes and Consani let \(H=x\partial_x\) be the scaling generator and define
\(\Delta=H(1+H)\).  On the quotient of the strong Schwartz space
\(\mathcal S_\infty(\mathbb R_+^*)\) by the closure of the range of the
trace map composed with \(\Delta\), they obtain the spectral values
\begin{equation}
\label{eq:spc-CC-laplacian-spectrum}
 \left(\rho-\frac12\right)^2-\frac14
 \quad
 (\zeta(\rho)=0,\ \rho\notin\mathbb R),
\end{equation}
counted with possible multiplicities
\cite[Prop.~4.2]{connesConsani2023Hochschild}.

Define on the spectral plane
\begin{equation}
\label{eq:spc-laplacian-coordinate}
 \lambda_\Delta(z)=-z^2-\frac14.
\end{equation}

\begin{proposition}[Exact Laplacian fibres and symmetry transport]
\label{prop:spc-laplacian-fibres}
For every \(s\in\mathbb C\),
\begin{equation}
\label{eq:spc-laplacian-affine-identity}
 \lambda_\Delta(\Phi(s))
 =
 -s(1-s)
 =
 \left(s-\frac12\right)^2-\frac14.
\end{equation}
For any \(z_1,z_2\in\mathbb C\),
\begin{equation}
\label{eq:spc-laplacian-fibre-equivalence}
 \lambda_\Delta(z_1)=\lambda_\Delta(z_2)
 \quad\Longleftrightarrow\quad
 z_2=z_1\ \hbox{or}\ z_2=-z_1.
\end{equation}
Thus the fibre is \(\{z,-z\}\) for \(z\ne0\) and is the singleton
\(\{0\}\) at the branch point.  The branch point is absent from
\(\Phi(|D_\xi|)\), since \(\xi(1/2)>0\).  Consequently, if
\(\operatorname{Spec}_{\mathrm{val}}(\Delta)\) denotes the underlying set
of distinct spectral values in \eqref{eq:spc-CC-laplacian-spectrum},
\begin{equation}
\label{eq:spc-sign-quotient-spectrum-bijection}
 \Phi(|D_\xi|)/\langle z\mapsto-z\rangle
 \overset{\sim}{\longrightarrow}
 \operatorname{Spec}_{\mathrm{val}}(\Delta),
 \qquad
 [z]\longmapsto\lambda_\Delta(z),
\end{equation}
is a bijection.  Its inverse sends a value \(\lambda\)
to the unordered root set of
\(z^2=-(\lambda+1/4)\).

The two packet involutions act by
\begin{equation}
\label{eq:spc-laplacian-involutions}
 \lambda_\Delta(-z)=\lambda_\Delta(z),
 \qquad
 \lambda_\Delta(-\overline z)=
 \overline{\lambda_\Delta(z)}.
\end{equation}
Hence, on \(\Phi(|D_\xi|)\), a full packet maps to the conjugation orbit
\(\{\lambda,\overline\lambda\}\), with repetitions removed.  This orbit
is a singleton precisely when \(z\in\mathbb R\), equivalently when the
original zero packet is critical and has two points.
\end{proposition}

\begin{proof}
Using \(\Phi(s)=\mathrm i(s-1/2)\),
\[
 -\Phi(s)^2-\frac14
 =
 \left(s-\frac12\right)^2-\frac14
 =s^2-s=-s(1-s),
\]
which proves \eqref{eq:spc-laplacian-affine-identity}.  Equality of two
values of \(\lambda_\Delta\) is equality of their squares:
\[
 z_1^2=z_2^2
 \quad\Longleftrightarrow\quad
 (z_1-z_2)(z_1+z_2)=0,
\]
proving the complete fibre formula.  The only collapsed fibre is at zero;
its inverse affine point is \(1/2\), and
\cref{prop:cpp-centered-symmetry} proves \(\xi(1/2)>0\).  Therefore no
exceptional fibre occurs on the completed-zero support, and the quotient map
is bijective with the stated root-set inverse.

The identities in \eqref{eq:spc-laplacian-involutions} are direct
substitutions.  By \cref{prop:spc-affine-strip}, the functional involution is
the sign map and conjugation is \(z\mapsto-\overline z\); therefore the image
of the four-point packet has the two displayed conjugate values.  If
\(z=x+\mathrm iy\), then
\[
 \Im\lambda_\Delta(z)=-2xy.
\]
On \(\Phi(|D_\xi|)\), the alternative \(x=0\) would make
\(\Phi^{-1}(z)\) a real point of \((0,1)\), where \(\zeta\) has no zero by
the eta-series argument in \cref{prop:cpp-centered-symmetry}.  Hence a real
spectral value on this support forces \(y=0\), which is exactly a real
spectral packet.  The converse is immediate.
\end{proof}

\subsection{A real polynomial coordinate for a whole packet}

For \(O\in Y_\xi\), let
\[
 \operatorname{Sp}_\Delta(O)
 =
 \{\lambda_\Delta(\Phi(\rho)):\rho\in O\}
\]
with repetitions removed, and define the monic spectral packet polynomial
\begin{equation}
\label{eq:spc-spectral-polynomial-definition}
 R_O(t)=\prod_{\lambda\in\operatorname{Sp}_\Delta(O)}(t-\lambda).
\end{equation}

\begin{theorem}[Spectral packet polynomial and coordinate inverse]
\label{thm:spc-spectral-packet-polynomial}
The polynomial \(R_O\) lies in \(\mathbb R[t]\), and the map
\begin{equation}
\label{eq:spc-spectral-polynomial-map}
 \mathcal R:Y_\xi\longrightarrow\mathbb R[t],
 \qquad O\longmapsto R_O,
\end{equation}
is injective.

If
\[
 O=\left\{\frac12+\mathrm i\gamma,
           \frac12-\mathrm i\gamma\right\},
 \qquad \gamma\in\mathbb R\setminus\{0\},
\]
then
\begin{equation}
\label{eq:spc-critical-spectral-polynomial}
 R_O(t)=t+\gamma^2+\frac14.
\end{equation}
If \(O\) has four points and
\[
 \alpha=a+\mathrm i\gamma,
 \qquad
 O=\left\{\frac12\pm\alpha,
           \frac12\pm\overline\alpha\right\},
 \qquad a\gamma\ne0,
\]
then
\begin{equation}
\label{eq:spc-noncritical-spectral-polynomial}
 R_O(t)
 =
 \left(t-\alpha^2+\frac14\right)
 \left(t-\overline{\alpha}^{2}+\frac14\right).
\end{equation}
The inverse takes the roots \(\lambda\) of \(R_O\), forms
\(A=\lambda+1/4\), takes the complete unordered square-root fibre
\(\{\alpha,-\alpha\}\) of \(A\), adjoins its conjugate fibre, and finally
translates every point by \(1/2\).  No square-root branch is selected.
\end{theorem}

\begin{proof}
The sign pair \(\{\rho,1-\rho\}\) has centered coordinates
\(\{\alpha,-\alpha\}\).  By
\eqref{eq:spc-laplacian-affine-identity}, both points give the same spectral
value \(\alpha^2-1/4\).  Conjugation supplies
\(\overline{\alpha}^{2}-1/4\).  This proves
\eqref{eq:spc-noncritical-spectral-polynomial}; its two factors are conjugate,
so their product has real coefficients.  In the critical case
\(\alpha=\mathrm i\gamma\), the two conjugate values coincide and equal
\(-\gamma^2-1/4\), giving \eqref{eq:spc-critical-spectral-polynomial}.

The roots of \(R_O\) recover the unordered spectral-value set.  Adding
\(1/4\) recovers the unordered set
\(\{\alpha^2,\overline{\alpha}^{2}\}\).  The complete square-root fibres are
\(\{\pm\alpha\}\) and \(\{\pm\overline\alpha\}\); their union is exactly
the centered packet, with repetitions removed in the critical case.
Translation by \(1/2\) therefore recovers \(O\).  This is an explicit inverse
to \(\mathcal R\), proving injectivity.
\end{proof}

The packet polynomial \(P_O\) of Chapter~17 and the spectral polynomial
\(R_O\) contain the same packet coordinate in different target spaces.  The
next result gives both directions at coefficient level.

\begin{proposition}[Exact conversion from packet coefficients to spectral coefficients]
\label{prop:spc-packet-spectral-coefficient-conversion}
Write the centered packet polynomial as
\[
 P_O\left(\frac12+u\right)=1+c_2u^2
\]
in the two-point case, or
\[
 P_O\left(\frac12+u\right)=1+c_2u^2+c_4u^4
\]
in the four-point case.

In the two-point case,
\begin{equation}
\label{eq:spc-critical-coefficient-conversion}
 R_O(t)=t-\lambda_O,
 \qquad
 \lambda_O=-\frac1{c_2}-\frac14.
\end{equation}
In the four-point case,
\begin{equation}
\label{eq:spc-quartic-coefficient-conversion}
 R_O(t)
 =
 t^2+\left(\frac{c_2}{c_4}+\frac12\right)t
 +\left(\frac1{c_4}+\frac{c_2}{4c_4}+\frac1{16}\right).
\end{equation}
Conversely, if \(\lambda\) is either root of the latter polynomial and
\(A=\lambda+1/4\), then the squared real coordinates of the packet are
\begin{equation}
\label{eq:spc-spectral-coordinate-inverse}
 a^2=\frac{|A|+\Re A}{2},
 \qquad
 \gamma^2=\frac{|A|-\Re A}{2}.
\end{equation}
The alternative root conjugates \(A\) and leaves these two numbers unchanged.
Their sign choices are exactly the four packet points, not discarded data.
Thus \(P_O\leftrightarrow R_O\) is a bijective change of packet coordinate.
\end{proposition}

\begin{proof}
In the critical case, \cref{prop:cpp-packet-polynomial-coordinates} gives
\(c_2=\gamma^{-2}\), while
\eqref{eq:spc-critical-spectral-polynomial} has root
\(\lambda_O=-\gamma^2-1/4\).  Substitution proves
\eqref{eq:spc-critical-coefficient-conversion}.

For a four-point packet, put \(A=\alpha^2\) and
\(B=\overline{\alpha}^{2}\).  The factorization from Chapter~17 is
\[
 P_O\left(\frac12+u\right)
 =\left(1-\frac{u^2}{A}\right)
  \left(1-\frac{u^2}{B}\right).
\]
Coefficient comparison gives
\begin{equation}
\label{eq:spc-AB-from-c}
 AB=\frac1{c_4},
 \qquad
 A+B=-\frac{c_2}{c_4}.
\end{equation}
The two spectral roots are \(\lambda=A-1/4\) and
\(\overline\lambda=B-1/4\).  Therefore
\[
 \lambda+\overline\lambda
 =-\frac{c_2}{c_4}-\frac12
\]
and
\[
 \lambda\overline\lambda
 =AB-\frac{A+B}{4}+\frac1{16}
 =\frac1{c_4}+\frac{c_2}{4c_4}+\frac1{16}.
\]
Expanding \((t-\lambda)(t-\overline\lambda)\) proves
\eqref{eq:spc-quartic-coefficient-conversion}.

For the inverse, write \(\alpha=a+\mathrm i\gamma\).  Then
\[
 A=\alpha^2=(a^2-\gamma^2)+2\mathrm ia\gamma,
 \qquad
 |A|=|\alpha|^2=a^2+\gamma^2.
\]
Adding and subtracting these last two real equations gives
\eqref{eq:spc-spectral-coordinate-inverse}.  Replacing \(A\) by
\(\overline A\) changes the sign of \(a\gamma\) but not \(a^2\) or
\(\gamma^2\).  The four choices
\((a,\gamma),(-a,-\gamma),(a,-\gamma),(-a,\gamma)\) are exactly
\(\alpha,-\alpha,\overline\alpha,-\overline\alpha\), so the complete packet
fibre is retained.  Since \cref{thm:cpp-packet-polynomial-inverse} recovers
\(P_O\) from that same packet, the two displayed constructions are mutual
coordinate conversions.
\end{proof}

\subsection{Divisor multiplicity under the quadratic spectral map}

Let \(D_\Phi=\Phi_*D_\xi\).  The polynomial
\(\lambda_\Delta(z)=-z^2-1/4\) is proper.  Explicitly, if a compact set
\(K\subset\mathbb C\) contains \(\lambda_\Delta(z)\), then
\[
 |z|^2=|z^2|=|\lambda_\Delta(z)+1/4|
 \leq \max_{\lambda\in K}|\lambda+1/4|,
\]
so \(\lambda_\Delta^{-1}(K)\) is closed and bounded.  Since \(\xi\) is a
nonzero entire function, \(|D_\xi|\) is a closed discrete zero set; hence
\(\Phi(|D_\xi|)\) is closed and discrete as well.  Its intersection with the
compact set \(\lambda_\Delta^{-1}(K)\) is therefore finite.  Thus the following
divisor pushforward is locally
finite, not merely finite on each point fibre:
\[
 (\lambda_\Delta)_*D_\Phi
 =
 \sum_\lambda
 \left(
   \sum_{z\in\Phi(|D_\xi|),\ \lambda_\Delta(z)=\lambda}m_z
 \right)[\lambda].
\]

\begin{corollary}[Exact aggregation of zero multiplicities]
\label{cor:spc-divisor-pushforward}
Let \(O\in Y_\xi\) have common multiplicity \(m_O\).
If \(O\) is critical, its contribution to
\((\lambda_\Delta)_*D_\Phi\) is
\begin{equation}
\label{eq:spc-critical-divisor-pushforward}
 2m_O[\lambda_O].
\end{equation}
If \(O\) has four points, its contribution is
\begin{equation}
\label{eq:spc-four-divisor-pushforward}
 2m_O[\lambda_O]+2m_O[\overline\lambda_O].
\end{equation}
No two distinct sign orbits contribute to the same scalar value.

This aggregation is compatible with the multiplicity mechanism proved by
Connes--Consani: a zero \(\rho\) of order \(m_O\) supplies an
\(m_O\)-dimensional jet space, and multiplication by \(-s(1-s)\) has a
triangular matrix with that spectral value repeated \(m_O\) times on the
diagonal \cite[proof of Prop.~4.2]{connesConsani2023Hochschild}.  Applying
their per-zero calculation to the two distinct zeros \(\rho\) and
\(1-\rho\), which have the same spectral value, gives two such jet blocks.
Their algebraic multiplicities therefore aggregate to the coefficient
\(2m_O\) in this chapter's divisor pushforward.
\end{corollary}

\begin{proof}
At every point of \(D_\Phi\), the affine map preserves the multiplicity
\(m_O\).  By \eqref{eq:spc-laplacian-fibre-equivalence}, a scalar spectral
value has exactly the sign fibre \(\{z,-z\}\), and both points occur in the
same packet with multiplicity \(m_O\).  Their sum is \(2m_O\).  A critical
packet consists of one sign fibre and hence gives
\eqref{eq:spc-critical-divisor-pushforward}.  A four-point packet is the
disjoint union of the sign fibres
\(\{z,-z\}\) and \(\{\overline z,-\overline z\}\), giving the two conjugate
terms in \eqref{eq:spc-four-divisor-pushforward}.  Finally, equality of two
spectral values forces equality up to sign by
\eqref{eq:spc-laplacian-fibre-equivalence}; hence distinct sign orbits cannot
collide.  The final paragraph is the direct coordinate comparison with the
cited jet calculation.
\end{proof}

\subsection{Critical packets: spectral value to every cycle length}

Let \(Y_\xi^{\mathrm{crit}}\) be the set of two-point packets.  For
\(O\in Y_\xi^{\mathrm{crit}}\), let \(\lambda_O\) be the unique root of
\(R_O\).  Then \(\lambda_O<-1/4\), and define
\begin{equation}
\label{eq:spc-critical-frequency-from-lambda}
 \nu(O)=\sqrt{-\lambda_O-\frac14}>0.
\end{equation}

\begin{proposition}[The cycle family factors through the Laplacian value]
\label{prop:spc-critical-cycle-factor}
For every \(O\in Y_\xi^{\mathrm{crit}}\) and every
\(n\in\mathbb Z_{\geq1}\), set
\begin{equation}
\label{eq:spc-cycle-family-from-lambda}
 \ell_n(O)
 =
 \frac{2\pi n}{\nu(O)}
 =
 \frac{2\pi n}{\sqrt{-\lambda_O-1/4}}.
\end{equation}
Then \(\ell_n(O)\) is exactly the Connes--Consani circle-length family
attached to the positive critical ordinate represented by \(O\).  The map
factors as
\[
 O\longmapsto R_O
 \longmapsto\lambda_O
 \longmapsto\nu(O)
 \longmapsto(\ell_n(O))_{n\geq1},
\]
and every arrow has the explicit inverse or fibre described above.  For this
fixed packet, \(\ell_1(O)\) is the least member of the displayed
positive-integer family.  This is not a comparison with lengths arising from
other ordinates.
\end{proposition}

\begin{proof}
Write
\(O=\{1/2+\mathrm i\gamma,1/2-\mathrm i\gamma\}\) with
\(\gamma\ne0\).  Equation
\eqref{eq:spc-critical-spectral-polynomial} gives
\[
 \lambda_O=-\gamma^2-\frac14,
 \qquad
 \nu(O)=|\gamma|.
\]
Thus \eqref{eq:spc-cycle-family-from-lambda} is
\(2\pi n/|\gamma|\).  Connes and Consani prove that if
\(s=|\gamma|>0\) and \(\zeta(1/2+\mathrm is)=0\), every circle whose length
is an integral multiple of \(2\pi/s\) is a zeta-cycle
\cite[Thm.~6.4(ii)]{connesConsani2021ZetaCycles}.  This is precisely the
displayed family.  Since \(n\) ranges over positive integers, its least member
for this fixed \(s\) occurs at \(n=1\); no comparison across different
values of \(s\) is involved.
\end{proof}

\subsection{What the spectral polynomial remembers about a leading jet}

For \(\rho\in O\), let
\[
 \mu_\rho=\frac{\xi^{(m_O)}(\rho)}{m_O!},
 \qquad
 \omega(O)=|\mu_\rho|.
\]
The transport calculation of \cref{prop:ds-phase-transport} proves that
\(\omega(O)\) is independent of the representative.

\begin{proposition}[Exact jet-weight dependence and one-packet information loss]
\label{prop:spc-jet-weight-information}
The spectral coordinate \(R_O\) determines \(P_O\), the packet \(O\), and
the common multiplicity once \(m_O\) is adjoined.  It consequently determines
\(|P_O'(\rho)|^{m_O}\), independently of the representative
\(\rho\in O\).  The actual packet weight satisfies
\begin{equation}
\label{eq:spc-weight-factorization}
 \omega(O)
 =
 |H_O(\rho)|\,|P_O'(\rho)|^{m_O},
\end{equation}
where
\[
 H_O(s)=\xi(1/2)
 \prod_{O'\ne O}P_{O'}(s)^{m_{O'}}.
\]
Thus a single \((R_O,m_O)\) does not determine \(\omega(O)\) in the class of
entire functions with the same centered sign and conjugation symmetries.  The
missing coordinate is exactly the nonzero complementary value
\(|H_O(\rho)|\).  The complete global packet system
\[
 \left(\xi(1/2),\{(R_{O'},m_{O'}):O'\in Y_\xi\}\right)
\]
does determine \(\omega(O)\), every leading jet, and the normalized local
germ.
\end{proposition}

\begin{proof}
The mutual conversion of \cref{prop:spc-packet-spectral-coefficient-conversion}
recovers \(P_O\) from \(R_O\).  Because
\(P_O(1/2+u)\) is even in \(u\) and has real coefficients,
\[
 |P_O'(\rho)|
 =|P_O'(1-\rho)|
 =|P_O'(\overline\rho)|
 =|P_O'(1-\overline\rho)|.
\]
Hence its absolute derivative is packet-valued.  The global-to-local theorem
\cref{thm:cpp-global-to-local-jet} proves the exact, unnormalized identity
\[
 \mu_\rho=H_O(\rho)(P_O'(\rho))^{m_O};
\]
taking absolute values gives \eqref{eq:spc-weight-factorization}.

To see that the complementary factor cannot be inferred from one packet,
retain the explicit fibre witness of
\cref{prop:cpp-one-packet-information-loss}.  For a positive real
\(\lambda\) chosen so that
\(Q_\lambda(s)=1+\lambda(s-1/2)^2\) has no root in \(O\) and
\(Q_\lambda(\rho)\ne1\), the functions
\[
 P_O(s)^{m_O}
 \quad\hbox{and}\quad
 Q_\lambda(s)P_O(s)^{m_O}
\]
  have the same local packet polynomial and multiplicity at \(O\), but their
  leading coefficients differ by the nonzero, nonidentity scalar factor
  \(Q_\lambda(\rho)\ne0,1\).
Both retain evenness about \(1/2\) and conjugation symmetry.  This proves the
nontrivial fibre of the one-packet projection at the exact claimed scope.  If
all packet factors and \(\xi(1/2)\) are supplied, the displayed product gives
\(H_O\) and hence every quantity in the statement.
\end{proof}

\subsection{The Scaling Site quotient and the shared coordinate diagram}

The second published realization uses a different, fully specified sheaf.
For \(L>0\), put \(\mu=e^L\) and
\(C_\mu=\mathbb R_+^*/\mu^{\mathbb Z}\).  Connes and Consani construct a
sheaf \(\mathcal L^2\) whose fibre is \(L^2(C_\mu)\), a closed subsheaf
\(\overline{\Sigma\mathcal E}\), and prove the equivariant isomorphism
\begin{equation}
\label{eq:spc-scaling-site-isomorphism}
 H^0\left(\mathfrak S_{\mathrm{sc}},
          \mathcal L^2/\overline{\Sigma\mathcal E}\right)
 \ \cong\ 
 \frac{\mathcal S(\mathbb R)}
 {\overline{\zeta(1/2+\mathrm i\,\cdot)\mathcal S(\mathbb R)}}.
\end{equation}
Here the closure is in the Schwartz topology, and \(r\in\mathbb R_+^*\)
acts at the Schwartz coordinate \(s\) by multiplication with \(r^{\mathrm is}\)
\cite[the theorem and proof on pp.~96--98]{connesConsani2023Hochschild}.
The pointwise Fourier transform on \(C_{e^L}\) uses the kernel
\(u^{-2\pi\mathrm ik/L}\), and the published scaling action on its component
of index \(k\) is multiplication by \(r^{-2\pi\mathrm ik/L}\).  Therefore the
exact index-to-Schwartz-coordinate map is
\begin{equation}
\label{eq:spc-signed-frequency-coordinate}
 s_k(L)=-\frac{2\pi k}{L},
 \qquad
 \zeta\left(\frac12-\frac{2\pi\mathrm ik}{L}\right)
 =\zeta\left(\frac12+\mathrm i s_k(L)\right),
 \qquad
 \chi_k(r)=r^{\mathrm i s_k(L)}=r^{-2\pi\mathrm ik/L}.
\end{equation}
Thus the sign of the integer Fourier index is opposite to the sign of the
Schwartz spectral coordinate; neither sign is normalized away.  Their proof
identifies the \(\pm1\) Fourier components, proves surjectivity of the
global-section quotient map, and compares the two closed ideals through this
coordinate substitution.  Their following remark records the nilpotent Jordan
term that can occur at a multiple critical zero.

\begin{theorem}[Positive interface between a critical packet and the Scaling Site]
\label{thm:spc-scaling-site-interface}
Let \(O\in Y_\xi^{\mathrm{crit}}\), let
\(\nu=\sqrt{-\lambda_O-1/4}\), and let \(n\geq1\).  Put
\begin{equation}
\label{eq:spc-L-mu-frequency}
 L=\frac{2\pi n}{\nu},
 \qquad
 \mu=e^L.
\end{equation}
Then the absolute frequency of the Fourier components with indices
\(+n\) and \(-n\) in \(L^2(C_\mu)\) is
\begin{equation}
\label{eq:spc-Fourier-frequency-equality}
 \frac{2\pi n}{L}=\nu,
\end{equation}
and
\begin{equation}
\label{eq:spc-scaling-zeta-vanishing}
 \zeta\left(\frac12+\mathrm i\nu\right)
 =
 \zeta\left(\frac12-\mathrm i\nu\right)
 =0.
\end{equation}
The resulting circle is one of the zeta-cycles of
\cref{prop:spc-critical-cycle-factor}, and the Scaling Site spectral
characters of the two components are
\begin{equation}
\label{eq:spc-scaling-character}
 \chi_{+n}(r)=r^{-\mathrm i\nu},
 \qquad
 \chi_{-n}(r)=r^{\mathrm i\nu}.
\end{equation}
Equivalently,
\begin{equation}
\label{eq:spc-signed-index-frequency-pair}
 s_{+n}(L)=-\nu,
 \qquad
 s_{-n}(L)=+\nu.
\end{equation}
Accordingly, the following coordinate chain commutes by direct substitution:
\[
 \begin{array}{ccccc}
 O&\longmapsto&\lambda_O&\longmapsto&
 \nu=\sqrt{-\lambda_O-1/4}\\[2mm]
 &&&&\downarrow\scriptstyle n\\[-1mm]
 &&&&L=2\pi n/\nu\longmapsto C_{e^L},
 \end{array}
\]
while the Fourier coordinate returns \(2\pi n/L=\nu\).

More precisely, the displayed formulas define the map
\begin{equation}
\label{eq:spc-labeled-scaling-interface-map}
 \Theta(O,n)
 =
 \left(
   \lambda_O,\nu,L,C_{e^L};
   (+n,-\nu,r\mapsto r^{-\mathrm i\nu}),
   (-n,+\nu,r\mapsto r^{\mathrm i\nu})
 \right)
\end{equation}
from \(Y_\xi^{\mathrm{crit}}\times\mathbb Z_{\geq1}\) to the set of labeled
packet-value, positive-frequency, circle, Fourier-index, signed-frequency, and
unitary-character tuples.  This is the exact shared-coordinate morphism between
the packet, Laplacian, cycle, and Scaling Site presentations proved here.  It
does not select a nonzero vector or global section of the Scaling Site quotient,
and it does not identify the Boolean packet sheaf of
\cref{thm:ds-packet-descent} with
\(\mathcal L^2/\overline{\Sigma\mathcal E}\): a stalkwise map, its topology,
and its compatibility with the \(\mathbb N^\times\)-equivariant structure
have not been constructed in this chapter.
\end{theorem}

\begin{proof}
Choose the positive ordinate \(\nu\) represented by the critical packet.
By construction of \(Y_\xi^{\mathrm{crit}}\), the conjugate pair
\(1/2\pm\mathrm i\nu\) belongs to \(O\), and both points are zeros of
\(\zeta\).  This proves \eqref{eq:spc-scaling-zeta-vanishing}.
Substitution of \eqref{eq:spc-L-mu-frequency} gives
\[
 \frac{2\pi n}{L}
 =
 \frac{2\pi n}{2\pi n/\nu}
 =\nu,
\]
which is \eqref{eq:spc-Fourier-frequency-equality}.  In the Fourier formula
for the Scaling Site subsheaf, the factor for index \(+n\) is therefore
\[
 \zeta\left(\frac12-\frac{2\pi\mathrm in}{L}\right)
 =
 \zeta\left(\frac12-\mathrm i\nu\right)=0,
\]
one member of \eqref{eq:spc-scaling-zeta-vanishing}; index \(-n\) gives the
other member.  The published cycle theorem then supplies the zeta-cycle.
The Scaling Site action on the component of index \(k\) is
\(r^{-2\pi\mathrm ik/L}\); substituting \(k=+n\) and \(k=-n\) gives
exactly the two characters in \eqref{eq:spc-scaling-character}.  Every edge
of the diagram is one of the displayed formulas, so commutativity is literal.

The Boolean packet stalk \(A^{m_O}\) is a finite semimodule, whereas the
published Scaling Site object is the stated quotient of a sheaf of complex
topological vector spaces.  The positive coordinate chain just proved gives
a common base parameter, but no formula above defines a map between those
stalks.  The final sentence records precisely that remaining construction
obligation; it is not a claim that such a morphism cannot exist.
\end{proof}

\subsection{The resulting packet and the exact scope of the reconstruction}

\paragraph{Finite certification boundary.}
The companion Python certificate replays twenty named exact checks
in two Python/SymPy installations.  They cover the split-zero operation tables,
both affine inverses, the quadratic fibre, both spectral-polynomial coordinate
conversions, the squared-coordinate inverse, the two divisor aggregations, and
the cycle-frequency substitution.  The companion Lean module proves thirteen
named finite algebraic
theorems, including the complete sign-fibre equivalence, under Lean~4.32.0 and
the pinned Mathlib revision.  These certificates deliberately do not replace
the cited analytic inputs: analytic continuation, the common-zero theorem, the
Laplacian quotient spectrum, and the Scaling Site isomorphism remain the
human-authored results cited above.

Combining the preceding maps, a packet \(O\) with multiplicity \(m_O\) has
the following explicit data and no hidden representative convention:
\begin{equation}
\label{eq:spc-weighted-spectral-packet}
 \left(
   R_O(t),\quad
   m_O,\quad
   A_{\mathbb C}^{m_O},\quad
   \omega(O)
 \right).
\end{equation}
The polynomial \(R_O\) recovers the complete point packet; \(m_O\) recovers
its effective divisor; \(A_{\mathbb C}^{m_O}\) is the descended Boolean
multiplicity stalk; and \(\omega(O)\) is the representative-independent
absolute leading jet.  On the critical subspace, the unique root of \(R_O\)
also recovers every length in \eqref{eq:spc-cycle-family-from-lambda} and the
Scaling Site frequency in \eqref{eq:spc-Fourier-frequency-equality}.

Three source formulations have thereby been sharpened at their exact proof
sites.  First, the common-zero identification is the divisor equality on
\(|\Im z|<1/2\), not a global equality after the trivial zeros are restored.
Second, \(2\pi/\nu\) is the least length in the positive-integer family for a
fixed ordinate \(\nu\), not a globally shortest cycle claim.  Third,
\(G(\mathbb Z)\) and \(G(\mathbb C)\) are connected by the proved embedding
\eqref{eq:spc-semiring-embedding}; their different coefficient rings do not
erase the common Boolean fixed pair.  These are corrections and positive
morphisms relative to the audited local source.  No general claim about what
has or has not been proved elsewhere is inferred from a gap or wording defect
in that source.
